%! Author = jdubec
%! Date = 16/02/2026

\documentclass[11pt, a4paper]{article}

% ─── Packages ────────────────────────────────────────────────────────────────
\usepackage[utf8]{inputenc}
\usepackage[T1]{fontenc}
\usepackage[margin=2.5cm]{geometry}
\usepackage{graphicx}
\usepackage{hyperref}
\usepackage{booktabs}       % nicer tables (\toprule, \midrule, \bottomrule)
\usepackage{float}          % [H] placement for figures and tables
\usepackage{xcolor}         % colors (used by listings)
\usepackage{listings}       % code listings
\hypersetup{colorlinks=true, linkcolor=black, citecolor=black, urlcolor=blue!70!black}

% ─── Code listing style ─────────────────────────────────────────────────────
\lstset{
    basicstyle=\ttfamily\small,
    breaklines=true,
    frame=single,
    numbers=left,
    numberstyle=\tiny\color{gray},
    xleftmargin=2em,
    framexleftmargin=1.5em
}
% Pro tip: for syntax highlighting, look into the 'minted' package.

% ─── Bibliography ────────────────────────────────────────────────────────────
\usepackage[style=numeric, backend=bibtex]{biblatex}
\addbibresource{references.bib}

% ─── Language (swap to lang/sk for Slovak) ───────────────────────────────────
\input{lang/sk}

% ─── Student details ─────────────────────────────────────────────────────────
\newcommand{\StudentName}{Bizoň Jozef, Moško Marek}
\newcommand{\StudyGroup}{Rucek - Streda 8:00}
\newcommand{\AssignmentNumber}{1}
\newcommand{\AssignmentTitle}{Kino}
\newcommand{\AcademicYear}{2025/2026}

% =============================================================================
\begin{document}

% ─── Cover page ──────────────────────────────────────────────────────────────
\begin{titlepage}
    \centering

    \includegraphics[width=0.3\textwidth]{../assets/logo}

    \vspace{0.3cm}
    {\large \LabelUniversity}\\[2pt]
    {\large \LabelFaculty}

    \vspace{2.5cm}
    {\Huge\bfseries \LabelProtocol}

    \vspace{0.6cm}
    {\LARGE \LabelCode{} --- \LabelCourse}

    \vspace{0.4cm}
    {\Large \LabelAssignment{} \AssignmentNumber: \AssignmentTitle}

    \vfill

    \begin{tabular}{r l}
        \textbf{\LabelStudent:} & \StudentName  \\[4pt]
        \textbf{\LabelGroup:}   & \StudyGroup   \\[4pt]
        \textbf{\LabelYear:}    & \AcademicYear \\[4pt]
        \textbf{\LabelDate:}    & \today        \\
    \end{tabular}

    \vspace{1.5cm}
\end{titlepage}

% ─── Table of contents ───────────────────────────────────────────────────────
\renewcommand{\contentsname}{\LabelContents}
\tableofcontents
\newpage

% =============================================================================
% This document serves as both the assignment template and a quick guide
% to the LaTeX features you will need. Replace the content, keep the structure.
%
% Template source: https://github.com/FIIT-Databases/assignment-template
% =============================================================================

\section{Špecifikácia systému}
\label{sec:specifikacia}

Navrhovaný systém predstavuje informačný systém pre kino, ktorý umožňuje správu filmov,
premietaní, kinosál, zákazníkov a lístkov. Systém poskytuje funkcionality pre zákazníkov aj
zamestnancov kina.

Primárnym cieľom systému je umožniť zákazníkom jednoducho vyhľadávať filmy, kupovať
lístky na konkrétne premietania a sledovať históriu filmov, ktoré navštívili. Okrem toho
systém poskytuje možnosť ukladať filmy do zoznamu sledovania (watchlist), aby bol zákazník
informovaný o budúcich premietaniach.

Z pohľadu kina systém umožňuje spravovať kinosály, premietania filmov a evidovať
zamestnancov spolu s ich dochádzkou.

Systém tak centralizuje všetky informácie súvisiace s prevádzkou kina a zjednodušuje prácu
so zákazníkmi aj zamestnancami.

\section{Aktéri systému}
\label{sec:akteri}

\subsection{Zákazník}
Zákazník predstavuje osobu, ktorá využíva služby kina. Pomocou systému môže:
\begin{itemize}
    \item vyhľadávať filmy,
    \item prezerať premietania filmov,
    \item kupovať lístky,
    \item vracať lístky,
    \item ukladať filmy do watchlistu,
    \item prezerať históriu navštívených filmov.
\end{itemize}

Zákazník interaguje so systémom najmä pri rezervácii alebo nákupe lístkov.

\subsection{Zamestnanec}
Zamestnanec je pracovník kina, ktorý využíva systém na evidenciu pracovnej dochádzky.
Jeho hlavné úlohy v systéme sú:
\begin{itemize}
    \item zaznamenanie príchodu do práce,
    \item zaznamenanie odchodu z práce,
    \item evidencia pracovných záznamov.
\end{itemize}

\section{Doménový model systému}
\label{sec:domenovy-model}

Systém obsahuje viacero základných objektov (entít), ktoré reprezentujú reálne prvky kina.

\subsection{Kino}
Kino predstavuje organizáciu, ktorá prevádzkuje premietanie filmov. Kino obsahuje:
\begin{itemize}
    \item viacero kinosál,
    \item zamestnancov,
    \item plán premietaní.
\end{itemize}

\subsection{Kinosála}
Kinosála je miestnosť, v ktorej sa premietajú filmy. Každá kinosála má:
\begin{itemize}
    \item kapacitu miest,
    \item označenie.
\end{itemize}
Premietanie filmu sa vždy koná v konkrétnej kinosále.

\subsection{Film}
Film predstavuje audiovizuálne dielo, ktoré je premietané v kine. Film môže byť premietaný
viackrát a môže byť pridaný do watchlistu zákazníkov.

\subsection{Premietanie}
Premietanie predstavuje konkrétnu projekciu filmu v určitom čase a v konkrétnej kinosále.
Obsahuje informácie o:
\begin{itemize}
    \item filme,
    \item čase začiatku,
    \item cene lístka,
    \item kinosále.
\end{itemize}

\subsection{Lístok}
Lístok predstavuje rezerváciu alebo zakúpenie miesta na konkrétne premietanie filmu.
Obsahuje informácie o:
\begin{itemize}
    \item zákazníkovi,
    \item premietaní,
    \item mieste v sále,
    \item stave lístka,
    \item čase predaja.
\end{itemize}

\subsection{Miesto}

Miesto predstavuje konkrétne fyzické sedadlo v kinosále. Každé miesto je jednoznačne
identifikované v rámci sály kombináciou radu a čísla. Miesto obsahuje informácie o:

\begin{itemize}
    \item kinosále, v ktorej sa nachádza,
    \item rade sedadla,
    \item čísle sedadla,
    \item type miesta (napr. štandard, VIP).
\end{itemize}

Miesto slúži ako základ pre výber konkrétneho sedadla pri kúpe lístka. Každý lístok odkazuje
na práve jedno miesto, pričom to isté miesto nemôže byť predané dvakrát na to isté
premietanie.

\subsection{Watchlist}
Watchlist predstavuje zoznam filmov, ktoré si zákazník uložil, aby ich mohol sledovať.

\subsection{Zamestnanec}
Zamestnanec je osoba pracujúca v kine. Systém eviduje informácie o zamestnancoch a ich
pracovnej dochádzke.

\subsection{Dochádzka}
Dochádzka predstavuje evidenciu pracovného času zamestnanca a obsahuje:
\begin{itemize}
    \item čas príchodu,
    \item čas odchodu,
    \item poznámku.
\end{itemize}

\section{Biznis pravidlá systému}
\label{sec:biznis-pravidla}

Systém musí dodržiavať nasledujúce pravidlá:
\begin{enumerate}
    \item Premietanie musí byť priradené ku konkrétnemu filmu.
    \item Premietanie sa koná v konkrétnej kinosále.
    \item Zákazník môže vlastniť viacero lístkov.
    \item Lístok musí byť vždy priradený ku konkrétnemu premietaniu.
    \item Počet predaných lístkov nesmie prekročiť kapacitu kinosály.    
    \item Zákazník môže vrátiť lístok len pred začiatkom premietania.
    \item Film môže byť pridaný do watchlistu iba ak existuje v systéme.
    \item Zamestnanec môže mať viacero záznamov dochádzky.
\end{enumerate}

\section{Používateľské scenáre}
\label{sec:use-cases}

\subsection{UC1 -- Nákup lístka}
\textbf{Popis:} Zákazník si zakúpi lístok na konkrétne premietanie filmu.

\subsubsection*{Hlavný scenár}
\begin{enumerate}
    \item Zákazník vyhľadá film v systéme.
    \item Systém zobrazí dostupné premietania filmu.
    \item Zákazník si vyberie konkrétne premietanie.
    \item Systém zobrazí dostupné miesta v kinosále.
    \item Zákazník vyberie miesto.
    \item Systém overí dostupnosť miesta.
    \item Systém vytvorí lístok.
    \item Systém potvrdí nákup lístka.
\end{enumerate}

\subsubsection*{Alternatívny scenár A1 -- Miesto je už obsadené}
\begin{enumerate}
    \item Systém zistí, že vybrané miesto je už obsadené.
    \item Systém zobrazí upozornenie zákazníkovi.
    \item Zákazník si vyberie iné miesto.
\end{enumerate}

\subsubsection*{Dátová vrstva}
Systém vykoná \texttt{SELECT} dotaz do tabuliek \texttt{Premietanie}, \texttt{Kinosala} a
\texttt{Miesto}, aby overil existenciu premietania, prislúchajúcej sály a konkrétneho miesta.
Následne vykoná \texttt{SELECT} do tabuľky \texttt{Listok}, aby overil, či vybrané miesto nie je
už obsadené na dané premietanie. Ak je miesto dostupné, systém vykoná vloženie záznamu:
\begin{lstlisting}[language=SQL, caption={Vytvorenie lístka v UC1}]
INSERT INTO Listok (id_premietanie, id_zakaznik, id_miesto, cena, stav_listka, cas_predaja)
VALUES (...);
\end{lstlisting}

\begin{figure}[H]
    \centering
    \includegraphics[width=0.9\textwidth]{"../UCs/uc1.png"}
    \caption{Sekvenčný diagram -- UC1 Nákup lístka}
    \label{fig:uc1}
\end{figure}

\subsection{UC2 -- Pridanie filmu do watchlistu}

\subsubsection*{Hlavný scenár}
\begin{enumerate}
    \item Zákazník vyhľadá film.
    \item Zákazník zvolí možnosť \uv{Pridať do watchlistu}.
    \item Systém overí existenciu filmu.
    \item Systém uloží film do watchlistu.
\end{enumerate}

\subsubsection*{Alternatívny scenár A1 -- Film je už vo watchliste}
\begin{enumerate}
    \item Systém zistí, že film je už uložený vo watchliste zákazníka.
    \item Systém zobrazí upozornenie, že film už je sledovaný.
\end{enumerate}

\subsubsection*{Dátová vrstva}
Systém vykoná dotaz:
\begin{lstlisting}[language=SQL]
SELECT * FROM Film
WHERE id_film = ...;
\end{lstlisting}
aby overil existenciu filmu. Následne vykoná kontrolu v tabuľke \texttt{Watchlist}, či už záznam
existuje. Ak neexistuje, vykoná:
\begin{lstlisting}[language=SQL, caption={Pridanie filmu do watchlistu v UC2}]
INSERT INTO Watchlist (id_zakaznik, id_film, datum_pridania)
VALUES (...);
\end{lstlisting}

\begin{figure}[H]
    \centering
    \includegraphics[width=0.9\textwidth]{"../UCs/uc2.png"}
    \caption{Sekvenčný diagram -- UC2 Pridanie filmu do watchlistu}
    \label{fig:uc2}
\end{figure}

\subsection{UC3 -- Vrátenie lístka}

\subsubsection*{Hlavný scenár}
\begin{enumerate}
    \item Zákazník otvorí zoznam svojich lístkov.
    \item Vyberie lístok.
    \item Zvolí možnosť vrátenia lístka.
    \item Systém overí čas začiatku premietania.
    \item Systém označí lístok ako vrátený.
\end{enumerate}

\subsubsection*{Alternatívny scenár A1 -- Premietanie už začalo}
\begin{enumerate}
    \item Systém zistí, že aktuálny čas je po začiatku premietania.
    \item Systém vrátenie lístka neumožní.
\end{enumerate}

\subsubsection*{Dátová vrstva}
Systém vykoná \texttt{SELECT} z tabuliek \texttt{Listok} a \texttt{Premietanie}. Ak je čas premietania
ešte v budúcnosti, vykoná:
\begin{lstlisting}[language=SQL, caption={Označenie lístka ako vráteného v UC3}]
UPDATE Listok
SET stav_listka = 'vrateny'
WHERE id_listok = ...;
\end{lstlisting}

\begin{figure}[H]
    \centering
    \includegraphics[width=0.9\textwidth]{"../UCs/uc3.png"}
    \caption{Sekvenčný diagram -- UC3 Vrátenie lístka}
    \label{fig:uc3}
\end{figure}

\subsection{UC4 -- Zobrazenie histórie filmov}

\subsubsection*{Hlavný scenár}
\begin{enumerate}
    \item Zákazník otvorí svoj profil.
    \item Zvolí možnosť \uv{História filmov}.
    \item Systém vyhľadá lístky zákazníka.
    \item Systém zobrazí zoznam filmov.
\end{enumerate}

\subsubsection*{Alternatívny scenár A1 -- Zákazník ešte nenavštívil žiadne premietanie}
Systém zobrazí informáciu, že história je prázdna.

\subsubsection*{Dátová vrstva}
Systém vykoná \texttt{JOIN} medzi tabuľkami \texttt{Listok}, \texttt{Premietanie} a \texttt{Film}, aby
získal zoznam filmov navštívených zákazníkom.

\begin{figure}[H]
    \centering
    \includegraphics[width=0.9\textwidth]{"../UCs/uc4.png"}
    \caption{Sekvenčný diagram -- UC4 Zobrazenie histórie filmov}
    \label{fig:uc4}
\end{figure}

\subsection{UC5 -- Zaznamenanie dochádzky}

\subsubsection*{Hlavný scenár}
\begin{enumerate}
    \item Zamestnanec sa prihlási do systému.
    \item Zvolí možnosť \uv{Zaznamenať príchod}.
    \item Systém uloží čas príchodu.
\end{enumerate}

\subsubsection*{Alternatívny scenár A1 -- Zamestnanec zaznamenáva odchod}
\begin{enumerate}
    \item Zamestnanec zvolí možnosť \uv{Zaznamenať odchod}.
    \item Systém aktualizuje existujúci záznam dochádzky.
\end{enumerate}

\subsubsection*{Dátová vrstva}
Pri príchode systém vykoná:
\begin{lstlisting}[language=SQL]
INSERT INTO Dochadzka (id_zamestnanec, cas_prichodu, poznamka)
VALUES (...);
\end{lstlisting}
Pri odchode vykoná:
\begin{lstlisting}[language=SQL, caption={Aktualizácia dochádzky v UC5}]
UPDATE Dochadzka
SET cas_odchodu = ...
WHERE id_dochadzka = ...;
\end{lstlisting}

\begin{figure}[H]
    \centering
    \includegraphics[width=0.9\textwidth]{"../UCs/uc5.png"}
    \caption{Sekvenčný diagram -- UC5 Zaznamenanie dochádzky}
    \label{fig:uc5}
\end{figure}

\subsection{UC6 -- Zobrazenie histórie premietaní}

\subsubsection*{Hlavný scenár}
\begin{enumerate}
    \item Zákazník otvorí sekciu histórie.
    \item Systém vyhľadá premietania.
    \item Zobrazí zoznam filmov.
\end{enumerate}

\subsubsection*{Alternatívny scenár A1 -- Neexistujú žiadne premietania}
Systém zobrazí informáciu o prázdnom zozname.

\subsubsection*{Dátová vrstva}
Systém vykoná \texttt{SELECT} s \texttt{JOIN} medzi tabuľkami \texttt{Premietanie} a \texttt{Film}.

\begin{figure}[H]
    \centering
    \includegraphics[width=0.9\textwidth]{"../UCs/uc6.png"}
    \caption{Sekvenčný diagram -- UC6 Zobrazenie histórie premietaní}
    \label{fig:uc6}
\end{figure}

\section{Databázový model}
\label{sec:db-model}

Databázový model systému pozostáva z desiatich tabuliek:
\begin{itemize}
    \item \texttt{Kino},
    \item \texttt{Kinosala},
    \item \texttt{Miesto},
    \item \texttt{Film},
    \item \texttt{Premietanie},
    \item \texttt{Zakaznik},
    \item \texttt{Listok},
    \item \texttt{Watchlist},
    \item \texttt{Zamestnanec},
    \item \texttt{Dochadzka}.
\end{itemize}

Diagram databázy je uvedený na Obrázku~\ref{fig:er-diagram}.

\begin{figure}[H]
    \centering
    \includegraphics[width=\textwidth]{"../UCs/database ER diagram.png"}
    \caption{ER diagram databázy kina}
    \label{fig:er-diagram}
\end{figure}

\section{Technický popis tabuliek}
\label{sec:tabulky}

\subsection{Tabuľka \texttt{Kino}}
Účel tabuľky je evidovať základné informácie o kinách v systéme. Každé kino predstavuje
prevádzku, ktorá obsahuje viacero kinosál a zamestnáva pracovníkov.
\begin{description}
    \item[\texttt{id\_kino}] \texttt{int}, primárny kľúč, \texttt{IDENTITY}.
    \item[\texttt{nazov}] \texttt{varchar(100)}, \texttt{NOT NULL}.
    \item[\texttt{adresa}] \texttt{varchar(255)}, nepovinný atribút.
    \item[\texttt{kontakt}] \texttt{varchar(50)}, \texttt{UNIQUE}.
\end{description}
Primárny kľúč \texttt{id\_kino} je identifikátor. Vďaka tomu je možné kino jednoznačne identifikovať aj v prípade, že sa jeho údaje v budúcnosti zmenia, napríklad názov alebo adresa. Tabuľka je referencovaná z tabuliek \texttt{Kinosala} a \texttt{Zamestnanec}.

\subsection{Tabuľka \texttt{Kinosala}}
Tabuľka eviduje jednotlivé kinosály patriace ku konkrétnemu kinu. Každá kinosála predstavuje
priestor, v ktorom sa uskutočňujú premietania filmov.

\begin{description}
    \item[\texttt{id\_sala}] \texttt{int}, primárny kľúč, \texttt{IDENTITY}.
    \item[\texttt{id\_kino}] \texttt{int}, cudzí kľúč na \texttt{Kino}, \texttt{NOT NULL}.
    \item[\texttt{oznacenie}] \texttt{varchar(20)}, \texttt{NOT NULL}.
    \item[\texttt{kapacita}] \texttt{int}, \texttt{NOT NULL}, \texttt{CHECK (kapacita > 0)}.
\end{description}

Primárny kľúč \texttt{id\_sala} predstavuje každú kinosálu.
Cudzí kľúč \texttt{id\_kino} určuje, ku ktorému kinu kinosála patrí.
Atribút \texttt{oznacenie} slúži na označenie sály v rámci kina a atribút \texttt{kapacita}
udáva maximálny počet miest v danej sále.

\subsection{Tabuľka \texttt{Film}}
Tabuľka eviduje filmy dostupné v systéme. Každý film môže byť premietaný viackrát a môže byť
zaradený do watchlistu viacerých zákazníkov.

\begin{description}
    \item[\texttt{id\_film}] \texttt{int}, primárny kľúč, \texttt{IDENTITY}.
    \item[\texttt{nazov}] \texttt{varchar(100)}, \texttt{NOT NULL}.
    \item[\texttt{zaner}] \texttt{varchar(50)}, nepovinný atribút.
    \item[\texttt{datum\_premiery}] \texttt{date}, nepovinný atribút.
    \item[\texttt{dlzka\_minut}] \texttt{int}, \texttt{CHECK (dlzka\_minut > 0)}.
\end{description}

Primárny kľúč \texttt{id\_film} je identifikátor filmu.
Identifikátor filmu zostáva rovnaký aj v prípade, že sa zmenia jeho údaje, napríklad názov. Tabuľka je referencovaná z tabuliek \texttt{Premietanie} a \texttt{Watchlist}.

\subsection{Tabuľka \texttt{Premietanie}}
Tabuľka eviduje konkrétne premietania filmov v jednotlivých kinosálach. Každý záznam
reprezentuje jedno premietanie filmu v určitom čase.
\begin{description}
    \item[\texttt{id\_premietanie}] \texttt{int}, primárny kľúč.
    \item[\texttt{id\_film}] \texttt{int}, cudzí kľúč na \texttt{Film}, \texttt{NOT NULL}.
    \item[\texttt{id\_sala}] \texttt{int}, cudzí kľúč na \texttt{Kinosala}, \texttt{NOT NULL}.
    \item[\texttt{cas\_zaciatku}] \texttt{datetime}, \texttt{NOT NULL}.
    \item[\texttt{cena\_zakladna}] \texttt{decimal(10,2)}, \texttt{CHECK (cena\_zakladna > 0)}.
\end{description}
Primárny kľúč \texttt{id\_premietanie} predstavuje každé premietanie.
Cudzí kľúč \texttt{id\_film} odkazuje na tabuľku \texttt{Film} a určuje, ktorý film sa premieta.
Cudzí kľúč \texttt{id\_sala} odkazuje na tabuľku \texttt{Kinosala}, čím určuje, v ktorej sále sa
premietanie koná.

\subsection{Tabuľka \texttt{Miesto}}
Tabuľka eviduje jednotlivé sedadlá v kinosálach. Každé miesto patrí do konkrétnej
kinosály a je identifikované kombináciou radu a čísla sedadla.

\begin{description}
    \item[\texttt{id\_miesto}] \texttt{int}, primárny kľúč, \texttt{IDENTITY}.
    \item[\texttt{id\_sala}] \texttt{int}, cudzí kľúč na \texttt{Kinosala}, \texttt{NOT NULL}.
    \item[\texttt{rad}] \texttt{int}, \texttt{NOT NULL}.
    \item[\texttt{cislo\_miesta}] \texttt{int}, \texttt{NOT NULL}.
    \item[\texttt{typ\_miesta}] \texttt{varchar(20)}, \texttt{DEFAULT 'standard'}.
\end{description}

Primárny kľúč \texttt{id\_miesto} predstavuje každé miesto.
Cudzí kľúč \texttt{id\_sala} určuje, do ktorej kinosály miesto patrí.

Aby sa zabránilo duplicitným miestam v rámci jednej sály, je definované
obmedzenie \texttt{UNIQUE (id\_sala, rad, cislo)}.

\subsection{Tabuľka \texttt{Zakaznik}}
Tabuľka eviduje zákazníkov, ktorí využívajú služby kina. Zákazník môže nakupovať lístky,
vracať ich, prezerať históriu navštívených filmov a vytvárať si watchlist.

\begin{description}
    \item[\texttt{id\_zakaznik}] \texttt{int}, primárny kľúč, \texttt{IDENTITY}.
    \item[\texttt{meno}] \texttt{varchar(50)}, \texttt{NOT NULL}.
    \item[\texttt{priezvisko}] \texttt{varchar(50)}, \texttt{NOT NULL}.
    \item[\texttt{email}] \texttt{varchar(100)}, \texttt{NOT NULL}, \texttt{UNIQUE}.
    \item[\texttt{datum\_registracie}] \texttt{datetime}, \texttt{DEFAULT CURRENT\_TIMESTAMP}.
\end{description}

Primárny kľúč \texttt{id\_zakaznik} predstavuje každého zákazníka.
Atribút \texttt{email} je jedinečný, aby sa zabránilo duplicitnej registrácii rovnakého zákazníka.
Tabuľka je referencovaná z tabuliek \texttt{Listok} a \texttt{Watchlist}.

\subsection{Tabuľka \texttt{Listok}}
Tabuľka slúži na evidenciu predaných lístkov na jednotlivé premietania. Každý lístok
reprezentuje jedno konkrétne miesto na premietaní.
\begin{description}
    \item[\texttt{id\_listok}] \texttt{int}, primárny kľúč.
    \item[\texttt{id\_premietanie}] \texttt{int}, cudzí kľúč na \texttt{Premietanie}, \texttt{NOT NULL}.
    \item[\texttt{id\_zakaznik}] \texttt{int}, cudzí kľúč na \texttt{Zakaznik}, \texttt{NOT NULL}.
    \item[\texttt{id\_miesto}] \texttt{int}, cudzí kľúč na \texttt{Miesto}, \texttt{NOT NULL}.
    \item[\texttt{cena}] \texttt{decimal(10,2)}, \texttt{NOT NULL}.
    \item[\texttt{stav\_listka}] \texttt{varchar(20)}, \texttt{DEFAULT 'platny'}.
    \item[\texttt{cas\_predaja}] \texttt{datetime}, \texttt{DEFAULT CURRENT\_TIMESTAMP}.
\end{description}
Primárny kľúč \texttt{id\_listok} určuje každý predaný lístok. Cudzí kľúč
\texttt{id\_premietanie} určuje premietanie, na ktoré bol lístok zakúpený. Cudzí kľúč
\texttt{id\_zakaznik} odkazuje na zákazníka, ktorý lístok vlastní.
Cudzí kľúč \texttt{id\_miesto} určuje konkrétne miesto v kinosále, na ktoré bol lístok vydaný.
Atribút \texttt{cena} uchováva konečnú cenu lístka zaplatenú zákazníkom pri nákupe.
Aby nebolo možné predať rovnaké miesto viackrát na to isté premietanie, je nad atribútmi
\texttt{id\_premietanie} a \texttt{id\_miesto} definované obmedzenie
\texttt{UNIQUE (id\_premietanie, id\_miesto)}.

\subsection{Tabuľka \texttt{Watchlist}}
Tabuľka reprezentuje zoznam filmov, ktoré si zákazník uložil na sledovanie.
\begin{description}
    \item[\texttt{id\_watchlist}] \texttt{int}, primárny kľúč.
    \item[\texttt{id\_zakaznik}] \texttt{int}, cudzí kľúč na \texttt{Zakaznik}, \texttt{NOT NULL}.
    \item[\texttt{id\_film}] \texttt{int}, cudzí kľúč na \texttt{Film}, \texttt{NOT NULL}.
    \item[\texttt{datum\_pridania}] \texttt{datetime}, \texttt{NOT NULL}.
\end{description}
Primárny kľúč \texttt{id\_watchlist} určuje každý záznam vo watchliste. Tabuľka realizuje
vzťah M:N medzi zákazníkom a filmom, preto obsahuje cudzie kľúče \texttt{id\_zakaznik} a
\texttt{id\_film}. Aby nebolo možné pridať rovnaký film do watchlistu toho istého zákazníka
viackrát, je nad atribútmi \texttt{id\_zakaznik} a \texttt{id\_film} definované obmedzenie
\texttt{UNIQUE (id\_zakaznik, id\_film)}.

\subsection{Tabuľka \texttt{Zamestnanec}}
Tabuľka eviduje zamestnancov pracujúcich v kine.
\begin{description}
    \item[\texttt{id\_zamestnanec}] \texttt{int}, primárny kľúč.
    \item[\texttt{id\_kino}] \texttt{int}, cudzí kľúč na \texttt{Kino}, \texttt{NOT NULL}.
    \item[\texttt{meno}] \texttt{varchar(50)}, \texttt{NOT NULL}.
    \item[\texttt{priezvisko}] \texttt{varchar(50)}, \texttt{NOT NULL}.
    \item[\texttt{email}] \texttt{varchar(100)}, \texttt{UNIQUE}.
    \item[\texttt{pozicia}] \texttt{varchar(50)}, nepovinný atribút.
\end{description}
Primárny kľúč \texttt{id\_zamestnanec} id zamestnanca.
Cudzí kľúč \texttt{id\_kino} určuje, v ktorom kine zamestnanec pracuje.

\subsection{Tabuľka \texttt{Dochadzka}}
Tabuľka eviduje pracovnú dochádzku zamestnancov.
\begin{description}
    \item[\texttt{id\_dochadzka}] \texttt{int}, primárny kľúč.
    \item[\texttt{id\_zamestnanec}] \texttt{int}, cudzí kľúč na \texttt{Zamestnanec}, \texttt{NOT NULL}.    
    \item[\texttt{cas\_prichodu}] \texttt{datetime}, \texttt{NOT NULL}.
    \item[\texttt{cas\_odchodu}] \texttt{datetime}, nepovinný atribút.
    \item[\texttt{poznamka}] \texttt{varchar(255)}, nepovinný atribút.
\end{description}
Primárny kľúč \texttt{id\_dochadzka} predstavuje každý záznam dochádzky.
Cudzí kľúč \texttt{id\_zamestnanec} zabezpečuje prepojenie dochádzky so zamestnancom.

\section{Vzťahy a referenčná integrita}
\label{sec:vztahy}

\subsection{Kino -- Kinosala}
Typ vzťahu: 1:N. Jedno kino môže obsahovať viacero kinosál. Každá kinosála musí patriť
práve jednému kinu. Pri pokuse o vymazanie kina, ktoré obsahuje kinosály, je nastavené
pravidlo \texttt{ON DELETE RESTRICT}, aby sa zabránilo odstráneniu kina, ktoré má existujúce
sály.

\subsection{Kinosala -- Miesto}
Typ vzťahu: 1:N. Jedna kinosála obsahuje viacero miest. Každé miesto musí patriť
práve jednej kinosále. Tento vzťah je realizovaný pomocou cudzieho kľúča
\texttt{id\_sala} v tabuľke \texttt{Miesto}. Pri pokuse o vymazanie kinosály,
ku ktorej existujú evidované miesta, je vhodné nastaviť pravidlo
\texttt{ON DELETE RESTRICT}.

\subsection{Film -- Premietanie}
Typ vzťahu: 1:N. Jeden film môže mať viacero premietaní. Každé premietanie musí mať
priradený práve jeden film (povinný vzťah). Pri pokuse o vymazanie filmu, ktorý má
naplánované premietania, je nastavené pravidlo \texttt{ON DELETE RESTRICT}.

\subsection{Premietanie -- Listok}
Typ vzťahu: 1:N. Jedno premietanie môže mať viacero predaných lístkov. Pri zrušení
premietania je možné použiť pravidlo \texttt{ON DELETE CASCADE}, čo zabezpečí automatické
odstránenie všetkých lístkov na dané premietanie.

\subsection{Miesto -- Listok}
Typ vzťahu: 1:N. Jedno miesto môže byť použité v rôznych lístkoch pre rôzne
premietania. Na jednom konkrétnom premietaní však môže byť dané miesto
obsadené najviac raz. Tento vzťah je realizovaný pomocou cudzieho kľúča
\texttt{id\_miesto} v tabuľke \texttt{Listok}. Jedinečnosť obsadenia miesta
na konkrétnom premietaní zabezpečuje obmedzenie
\texttt{UNIQUE (id\_premietanie, id\_miesto)}.

\subsection{Zakaznik -- Listok}
Typ vzťahu: 1:N. Jeden zákazník môže vlastniť viacero lístkov. Každý lístok musí patriť
práve jednému zákazníkovi.

\subsection{Film -- Watchlist -- Zakaznik}
Typ vzťahu: M:N. Jeden zákazník môže sledovať viacero filmov a jeden film môže byť
sledovaný viacerými zákazníkmi. Tento vzťah je realizovaný pomocou tabuľky
\texttt{Watchlist}.

\subsection{Zamestnanec -- Dochadzka}
Typ vzťahu: 1:N. Jeden zamestnanec môže mať viacero záznamov dochádzky. Pri vymazaní
zamestnanca sa jeho záznamy dochádzky odstránia pomocou pravidla \texttt{ON DELETE CASCADE}.

\section{Zadanie 3}
\label{sec:zadanie3}

V rámci Zadania 3 implementujeme vybrané používateľské scenáre čistým SQL v PostgreSQL a podporujeme ich indexmi.
Nižšie uvádzame prvý z dvoch plánovaných procesov (druhý doplníme neskôr) a súhrn indexov zo súboru \texttt{indexy.sql}.

\subsection{Zmena oproti UC6 v Zadaní 1}

Use case \textbf{UC6 -- Zobrazenie histórie premietaní} bol v špecifikácii popísaný ako všeobecný výpis premietaní spojených s filmami (\texttt{JOIN} \texttt{Premietanie} -- \texttt{Film}) bez väzby na konkrétneho používateľa.
Pre potreby Zadania 3 sme proces upravili tak, aby zodpovedal realite profilu zákazníka: história sa zobrazuje \textbf{iba pre prihláseného zákazníka} a vychádza z jeho zakúpených lístkov (\texttt{Listok}), nie len zo zoznamu všetkých premietaní.
Zákazník tak vidí filmy (premietania), ktoré skutočne navštívil alebo si na ne zakúpil lístok (podľa filtrovania podľa stavu lístka).

\subsection{Proces 1 -- História premietaní z pohľadu zákazníka}

\textbf{Viazba na scenár:} nadväzuje na UC6 po úprave -- zákazník v profile vyžiada históriu svojich návštev súvisiacich s premietaniami.

\textbf{Vstup:} identifikátor zákazníka (\texttt{id\_zakaznik}), v ukážkovom skripte konštanta \texttt{42}.
Voliteľne možno sprísniť výber len na platné lístky (\texttt{stav\_listka = 'platny'}), aby sa do histórie nedostali vrátené alebo neplatné záznamy.

\textbf{Výstup:} pre každý relevantný lístok riadok s filmom, kinom, označením sály, časom začiatku premietania, miestom v sále (rad, číslo, typ), cenou a stavom lístka.
Dotaz navyše používa oknové funkcie (\texttt{SUM \ldots OVER}, \texttt{COUNT \ldots OVER}) na kumulatívnu útratu podľa času predaja, celkovú útratu a počet lístkov v rámci výslednej množiny -- čo predstavuje analytický prvok nad transakčnými dátami.

\textbf{SQL implementácia} je v súbore \texttt{historia\_premietani.sql}; jej podstata je nasledujúca:

\begin{lstlisting}[language=SQL, caption={História premietaní pre konkrétneho zákazníka (PostgreSQL)}]
SELECT
    l.id_listok,
    f.nazov AS film,
    k.nazov AS kino,
    ks.oznacenie AS sala,
    p.cas_zaciatku,
    m.rad,
    m.cislo_miesta,
    m.typ_miesta,
    l.cena,
    l.stav_listka,
    l.cas_predaja,
    SUM(l.cena) FILTER (WHERE l.stav_listka = 'platny')
        OVER (ORDER BY l.cas_predaja) AS kumulativne_utraty,
    SUM(l.cena) FILTER (WHERE l.stav_listka = 'platny')
        OVER () AS celkova_utrata,
    COUNT(*) OVER () AS pocet_listkov_celkom
FROM Listok l
JOIN Premietanie p ON p.id_premietanie = l.id_premietanie
JOIN Film f        ON f.id_film = p.id_film
JOIN Kinosala ks   ON ks.id_sala = p.id_sala
JOIN Kino k        ON k.id_kino = ks.id_kino
JOIN Miesto m      ON m.id_miesto = l.id_miesto
WHERE l.id_zakaznik = 42
  AND l.stav_listka = 'platny'
ORDER BY l.cas_predaja DESC;
\end{lstlisting}

\textbf{Meranie plánu:} v skripte sú poznámky z behu \texttt{EXPLAIN ANALYZE} pred a po pridaní indexov (orientačné časy plánovania a vykonania, počty bufferov).
Po nasadení kompozitného indexu na \texttt{(id\_zakaznik, stav\_listka)} klesol v našom meraní čas vykonania približne z $\sim$0{,}59\,ms na $\sim$0{,}27\,ms pri danom objeme dát.

\subsection{Druhý proces}

Druhý z dvoch procesov podľa zadania zatiaľ nie je v dokumentácii rozpracovaný; doplníme ho dodatočne spolu s vlastným dotazom a prípadnými ďalšími indexmi.

\subsection{Indexy (\texttt{indexy.sql})}

Pre dotaz na históriu lístkov zákazníka sú kľúčové podmienky \texttt{WHERE l.id\_zakaznik = \ldots} a filter podľa \texttt{stav\_listka}.
Definovali sme tieto indexy nad tabuľkou \texttt{Listok}:

\begin{description}
    \item[\texttt{idx\_listok\_zakaznik (id\_zakaznik)}]
        Podporuje vyhľadanie všetkých lístkov daného zákazníka bez nutnosti sekvenčného prechodu celej tabuľky \texttt{Listok}.
        Vyhýba sa nákladnému \emph{Seq Scan} pri výbere riadkov podľa zákazníka.

    \item[\texttt{idx\_listok\_stav (stav\_listka)}]
        Pomáha dotazom, ktoré filtrujú alebo agregujú podľa stavu lístka (\texttt{platny}, \texttt{vratený} a pod.).
        Môže byť užitočný aj v iných reportoch nad \texttt{Listok}.

    \item[\texttt{idx\_listok\_zakaznik\_stav (id\_zakaznik, stav\_listka)}]
        Kombinuje obe kritériá v jednom indexe v poradí zodpovedajúcom najčastejšiemu filtrovaniu v histórii profilu zákazníka.
        Plánovač môže efektívne zužovať množinu riadkov naraz podľa zákazníka aj stavu, čo znižuje počet čítaných blokov pred spojením s \texttt{Premietanie}, \texttt{Film}, \texttt{Kinosala}, \texttt{Kino} a \texttt{Miesto}.
\end{description}

Primárne a jedinečné kľúče (\texttt{id\_listok}, \texttt{UNIQUE (id\_premietanie, id\_miesto)}) už kryjú niektoré spojenia; uvedené indexy sú \emph{doplnkové} a zamerané konkrétne na proces histórie podľa zákazníka.

\section{Reflexia}
\label{sec:reflexia}

Pri návrhu databázového modelu sme postupovali postupne a viackrát sme ho upravovali podľa toho,
ako sme si spresňovali fungovanie celého systému. Na začiatku sme mali navrhnuté len základné
entity, ako napríklad \texttt{Film}, \texttt{Premietanie}, \texttt{Kinosala} a \texttt{Zakaznik}.
Neskôr sme dopĺňali ďalšie tabuľky, vzťahy medzi nimi a potrebné obmedzenia.

Najviac času nám zabrala časť okolo tabuliek \texttt{Film}, \texttt{Premietanie} a \texttt{Listok},
pretože bolo potrebné správne vyriešiť, ako evidovať konkrétne premietania filmu a zároveň aj
predaj miest na tieto premietania. Dôležité bolo najmä to, aby bolo možné priradiť lístok ku
konkrétnemu premietaniu a zároveň ku konkrétnemu miestu v sále.

Tému kina sme si vybrali preto, že je ľahko predstaviteľná a dobre sa na nej ukazujú vzťahy medzi
jednotlivými tabuľkami. Zároveň nám táto téma prišla dostatočne praktická na to, aby sme na nej
vedeli navrhnúť viacero zaujímavých scenárov a rôzne typy väzieb, napríklad vzťah 1:N medzi
filmom a premietaním alebo M:N medzi zákazníkom a filmom cez tabuľku \texttt{Watchlist}.

Pri kontrole dokumentu a overovaní, či si jednotlivé časti navzájom neodporujú, sme čiastočne
využili aj AI nástroje. Pomohli nám hlavne pri kontrole konzistencie medzi textovou časťou,
scenármi a databázovým modelom. Výstupy sme si však následne prešli, upravili a prispôsobili
podľa zadania.

% ─── Bibliografia (môže zostať prázdna) ───────────────────────────────────────
\printbibliography

\end{document}